\documentclass[
  aps,
  prx,
  reprint,
  amsmath,
  amssymb,
  groupedaddress,
  nofootinbib,
  floatfix,
  longbibliography
]{revtex4-2}

\usepackage[T1]{fontenc}
\usepackage{bm}
\usepackage{booktabs}
\usepackage{graphicx}
\usepackage{mathtools}
\usepackage{microtype}
\usepackage{xcolor}
\usepackage[colorlinks=true,allcolors=blue!55!black]{hyperref}

\setlength{\textfloatsep}{9pt plus 2pt minus 2pt}
\setlength{\floatsep}{8pt plus 2pt minus 2pt}
\setlength{\intextsep}{8pt plus 2pt minus 2pt}
\renewcommand{\arraystretch}{1.08}

\newcommand{\qefem}{QeFEM}
\newcommand{\lqa}{\textsc{LQA}}
\newcommand{\Tr}{\operatorname{Tr}}
\newcommand{\sign}{\operatorname{sign}}
\newcommand{\Dkl}{D_{\mathrm{KL}}}

\begin{document}

\title{A Quantum-Enhanced Free-Energy Machine for Large-Scale Binary Optimization}
\author{Ronit Dutta}
\affiliation{Singapore University of Technology and Design}
\author{Feng Pan}
\affiliation{Singapore University of Technology and Design}
\date{September 21, 2026}

\begin{abstract}
We present the Quantum-enhanced Free-Energy Machine (\qefem), a classical
variational method for large binary quadratic problems.  Each binary variable
is replaced by a one-qubit mixed state with longitudinal and transverse
polarizations.  Restricting a transverse-field Gibbs state to a product-state
family gives a closed-form free energy: temperature controls the entropy of
the relaxation, the transverse field supplies an additional continuous search
direction, and the final longitudinal polarization returns an ordinary binary
assignment.  Thousands of independently initialized product states can be
optimized in parallel with standard tensor operations.  \qefem{} reaches 33
of the stored reference cuts for G1--G54 without local postprocessing.  On the
dense 2000-vertex K2000 instance, a frozen five-trial evaluation attains cut
33,293, only 44 below the stored target 33,337, while local quantum annealing
(\lqa) attains 33,249 under its tuned configuration.  On fixed Wishart and
Chook planted suites, \qefem{} lowers the equal-instance mean relative energy
error from 6.977\% to 4.478\% and from 1.111\% to 0.2199\%, respectively.
These results show that a compact product-state relaxation, combined with
population search and careful conditioning, can return strong discrete
solutions on sparse, dense, and planted problem families.
\end{abstract}

\maketitle

\section{Introduction}
\label{sec:introduction}

Combinatorial optimization problems encode discrete choices in scheduling,
routing, allocation, circuit design, and scientific inference.  Many central
problems are NP-hard, so exact worst-case solution at scale is not expected
from a polynomial-time algorithm unless standard complexity assumptions fail
\cite{karp1972,barahona1982}.  Practical solvers therefore seek high-quality
feasible solutions while balancing generality, computational cost, and the
ability to exploit modern hardware.

A broad class of these problems can be written as quadratic unconstrained
binary optimization (QUBO), and hence as an Ising model
\cite{lucas2014}.  This correspondence connects combinatorial algorithms with
statistical mechanics and has motivated thermal annealing, coherent Ising
machines, digital annealers, and classically simulated nonlinear dynamics
\cite{kirkpatrick1983,inagaki2016science,aramon2019,goto2019sbm}.  Although
their dynamics differ, these methods share a basic strategy: replace immediate
enumeration of discrete assignments by a continuous or stochastic search that
can be evaluated efficiently.

The Free-Energy Machine (FEM) follows this strategy by minimizing a
differentiable mean-field free energy and using batched tensor operations to
explore many initial conditions in parallel \cite{shen2025}.  \qefem{}
develops the same principle in a mixed-state formulation.  Every variable is
represented by a one-qubit density matrix.  Its longitudinal polarization
encodes the eventual binary decision, its Bloch length determines the local
entropy, and its transverse polarization couples to a driver that is removed
during annealing.  A product-state approximation avoids the exponential cost
of the exact Gibbs state and leaves only two real parameters per variable and
replica.

The construction is quantum-inspired but the computation is classical.
\qefem{} does not require quantum hardware, sample from a quantum device, or
simulate entanglement.  Its value lies in the optimization geometry produced
by a transverse mixed-state relaxation.  At high temperature, entropy keeps
the local states weakly polarized.  As temperature and the transverse field decrease, the
problem couplings polarize the longitudinal components.  Final sign decoding
then produces a feasible binary solution.

This paper makes three contributions.  First, it derives the \qefem{}
functional from the binary objective through a transverse-field Gibbs target
and a Kullback--Leibler variational principle.  The derivation is kept in the
main text because it explains where every term in the optimizer comes from;
coefficient bookkeeping and analytic derivative checks are collected in the
appendices.  Second, it identifies the common algorithm beneath the original
automatic-differentiation implementation and the later nonlinear and
manual-gradient implementation.  They use the same local state family,
annealed free energy, replica population, and sign readout.  Third, it tests
the method on G-set, K2000, Wishart, and Chook instances, with \lqa{} as the
present comparison method \cite{bowles2022}.

The benchmark results are best understood in the scale of the underlying
problems.  A residual of 44 on K2000 is not merely a nonzero gap: it is a cut
of 33,293 on a complete 2000-vertex graph with 1,999,000 weighted edges.  On
G-set, \qefem{} reaches 33 of the first 54 stored references and leaves five
other instances a single cut unit short.  On the planted suites, it reduces
the mean relative error at every tested control parameter.  These outcomes
motivate the central question of the paper: how far can a simple factorized
free-energy geometry be pushed when it is carefully conditioned and explored
with a large replica population?  We answer that question by deriving the
relaxation first and fixing the evaluation protocol before interpreting the
results.

\section{From a binary objective to a free energy}
\label{sec:formulation}

\subsection{QUBO and Ising variables}

Let $\bm x\in\{0,1\}^{N}$, $N>1$ and let $Q$ be symmetric.  A QUBO objective and its
Ising representation are
\begin{equation}
 \begin{aligned}
  E_Q(\bm x)&=\bm x^{\mathsf T}Q\bm x+c,\\
  E_I(\bm s)&=c_0+\sum_i h_i s_i+\sum_{i<j}K_{ij}s_i s_j,
  \qquad s_i\in\{-1,+1\}.
 \end{aligned}
 \label{eq:qubo-ising}
\end{equation}
The substitution $x_i=(1-s_i)/2$ converts the first line into the second.
The constant $c_0$ changes objective values but not the minimizing spin
assignment.  Appendix~\ref{app:qubo} gives the coefficient map for the matrix
convention used here and records a standard quadratization penalty for
higher-order terms.

\subsection{MaxCut as an Ising energy}

For an undirected weighted graph $G=(V,E,w)$, the two spin values specify the
two sides of a cut.  An edge contributes its weight exactly when its endpoints
have opposite signs, so
\begin{equation}
 C(\bm s)=\frac12\sum_{\{i,j\}\in E}w_{ij}(1-s_i s_j).
 \label{eq:maxcut}
\end{equation}
Maximizing $C$ is equivalent to minimizing $E_P(\bm s)=-C(\bm s)$.  Replacing
each classical spin by the Pauli operator $Z_i$ produces a diagonal problem
Hamiltonian $H_P$ whose computational-basis ground states encode maximum cuts.
The same replacement converts a general Ising objective in
Eq.~\eqref{eq:qubo-ising} into a diagonal Hamiltonian.

\subsection{Why a Gibbs target appears}

Direct minimization still requires navigating $2^N$ assignments.  A Gibbs
state provides a controlled relaxation: finite temperature spreads weight
over many configurations, while the zero-temperature limit concentrates on
minimum-energy states.  \qefem{} augments the diagonal problem by a transverse
driver,
\begin{equation}
 H_{\Gamma}=H_P-\Gamma\sum_{i=1}^{N}X_i,
 \qquad
 \sigma_{T,\Gamma}=\frac{e^{-H_{\Gamma}/T}}
 {\Tr(e^{-H_{\Gamma}/T})}.
 \label{eq:gibbs-target}
\end{equation}
Here $T>0$ is temperature and $\Gamma\geq0$ controls the driver.  The driver
does not commute with the problem Hamiltonian and supplies a direction that is
orthogonal to the final computational-basis decision.  No dynamical quantum
evolution is simulated: at each optimization step, $T$ and $\Gamma$ simply
define the variational objective to be minimized.

The exact state in Eq.~\eqref{eq:gibbs-target} has dimension $2^N$ and is not
constructed.  Instead, choose a tractable trial state $\rho$.  Substituting
$\log\sigma_{T,\Gamma}=-H_{\Gamma}/T-\log Z$ into the quantum
Kullback--Leibler divergence gives
\begin{equation}
 \begin{aligned}
  \Dkl(\rho\Vert\sigma_{T,\Gamma})
  &=\Tr\!\left[\rho(\log\rho-\log\sigma_{T,\Gamma})\right]\\
  &=\frac{1}{T}\left[
      \Tr(\rho H_{\Gamma})-T S(\rho)-F^{\star}_{T,\Gamma}
    \right]\geq0,
 \end{aligned}
 \label{eq:kl-free-energy}
\end{equation}
where $S(\rho)=-\Tr(\rho\log\rho)$ and
$F^{\star}_{T,\Gamma}=-T\log Z$ is independent of the trial state.  Minimizing
the divergence is therefore exactly equivalent to minimizing
$\Tr(\rho H_{\Gamma})-TS(\rho)$.  This identity is the origin of the \qefem{}
free energy; the entropy term is not an unrelated regularizer
\cite{wainwright2008,blei2017,nielsen2010}.

\section{The QeFEM relaxation}
\label{sec:method}

\subsection{Product-state approximation}

The remaining difficulty is to choose a state family that is expressive
enough to smooth the binary landscape but inexpensive enough for thousands of
variables and replicas.  \qefem{} uses a fully factorized state,
\begin{equation}
 \begin{aligned}
  \rho&=\bigotimes_{i=1}^{N}\rho_i,\\
  \rho_i&=\frac12(I+m_{x,i}X_i+m_{z,i}Z_i),\\
  \binom{m_{x,i}}{m_{z,i}}
  &=\frac{\tanh r_i}{r_i}\binom{u_i}{v_i},
  \qquad r_i=(u_i^2+v_i^2)^{1/2}.
 \end{aligned}
 \label{eq:product-state}
\end{equation}
The unconstrained fields $(u_i,v_i)$ are mapped smoothly into the open Bloch
disk, so $m_{x,i}^2+m_{z,i}^2<1$ without projection or clipping.  The limit at
$r_i=0$ is well defined.  Measuring $Z_i$ gives probabilities
$q_i(s_i)=(1+s_i m_{z,i})/2$, and the product ansatz makes the joint
distribution $q(\bm s)=\prod_i q_i(s_i)$.  Consequently
$\langle Z_iZ_j\rangle=m_{z,i}m_{z,j}$ for $i\ne j$.  This factorization is the
mean-field approximation; the exact Gibbs state is not assumed to factorize.

The eigenvalues of $\rho_i$ are $(1\pm R_i)/2$, with
$R_i=(m_{x,i}^2+m_{z,i}^2)^{1/2}=\tanh r_i$.  Its entropy is therefore
\begin{equation}
 S_i=-\frac{1+R_i}{2}\log\frac{1+R_i}{2}
     -\frac{1-R_i}{2}\log\frac{1-R_i}{2}.
 \label{eq:local-entropy}
\end{equation}
At $R_i=0$ the site is maximally mixed; as $R_i\to1$ it becomes pure.  Since
entropy is additive on tensor products, $S(\rho)=\sum_iS_i$.

\subsection{Closed-form objective}

We can now evaluate every term in Eq.~\eqref{eq:kl-free-energy}.  For the Ising
problem in Eq.~\eqref{eq:qubo-ising}, product-state factorization replaces
$s_i$ by $m_{z,i}$ and $s_i s_j$ by $m_{z,i}m_{z,j}$.  The driver expectation
is $-\Gamma\sum_i m_{x,i}$.  Up to the irrelevant constant $c_0$, the complete
objective is therefore
\begin{equation}
 \boxed{\begin{aligned}
 \mathcal F_{T,\Gamma}
 &=\sum_i h_i m_{z,i}
  +\sum_{i<j}K_{ij}m_{z,i}m_{z,j}\\
 &\quad-\Gamma\sum_i m_{x,i}-T\sum_iS_i .
 \end{aligned}}
 \label{eq:qefem-objective}
\end{equation}
This equation contains the whole smooth \qefem{} model.  The first two terms
are the relaxed problem energy, the third is the transverse search direction,
and the fourth rewards uncertainty while the temperature is nonzero.  A new
QUBO changes only $(h_i,K_{ij})$; the state parameterization, entropy, driver,
and optimizer remain unchanged.

\subsection{Annealing, replicas, and readout}

At the start of a run, nonzero $T$ and $\Gamma$ keep the sites weakly polarized and allow
motion in both Bloch directions.  The schedules then lower $T$ and remove the
driver, so the graph-dependent longitudinal energy progressively dominates.
This is an optimization homotopy, not a claim of adiabatic quantum evolution.
The schedule family used in the experiments is stated in
Appendix~\ref{app:schedules}.

The free energy is nonconvex, and different initial fields can enter different
basins.  We therefore optimize $B$ independently initialized replicas in one
batched tensor.  RMSprop or SGD updates the unconstrained fields; momentum and
weight decay govern the numerical integration but do not alter
Eq.~\eqref{eq:qefem-objective}.  With an edge-list contraction, one update
costs $O(B|E|+BN)$ arithmetic and stores $O(BN+|E|)$ values.  Dense matrix
storage changes the leading arithmetic cost to $O(BN^2)$.

The returned solution is always discrete.  At the final update we set
$s_i=\sign(m_{z,i})$, rescore every replica with the original unscaled
couplings, and retain the best final assignment.  The continuous free energy
and expected cut are optimization diagnostics; all reported benchmark values
come from feasible binary states.  No local search or coordinate splicing is
applied.

\subsection{Equivalent implementations and the G-set continuation}

The early and current implementations look different but optimize the same
model.  Direct Bloch coordinates and the nonlinear fields in
Eq.~\eqref{eq:product-state} parameterize the same local mixed states.
Automatic differentiation and the exact manual update evaluate the same
smooth gradient of Eq.~\eqref{eq:qefem-objective}.  The nonlinear map improves
conditioning, while the manual form reduces overhead and makes sparse matrix
operations explicit; neither constitutes a different algorithm.  The
Jacobian and analytic gradients are given in Appendix~\ref{app:param}.

For difficult G-set instances, we additionally test a continuation in which
the neighbour magnetization entering the longitudinal search direction is
moved gradually toward its sign.  This changes the search direction but not
the state family, final readout, or reported energy, so it is identified
separately from the exact smooth gradient.  Its definition and activation
schedules are given in Appendix~\ref{app:schedules}.

The resulting distinction is precise: reparameterization and gradient
evaluation change how the same free energy is optimized, whereas the
discrete-neighbour continuation deliberately modifies the search direction.
With that distinction fixed, the next section states exactly which variant,
budget, and scoring rule is used for each benchmark.

\section{Experimental design}
\label{sec:experiments}

\subsection{Benchmark families and metrics}

The experiments cover four complementary settings.  G1--G54 include random,
toroidal, and planar MaxCut graphs with 800--3000 vertices \cite{yeGset}.
K2000 is a complete signed graph with 2000 vertices and 1,999,000 weighted
edges \cite{inagaki2016science}.  Wishart contains ten fixed 500-spin planted
instances indexed by $\alpha=0.1,\ldots,1.0$ \cite{hamze2020}.  Chook contains
eleven fixed periodic $32\times32$ tile-planted instances indexed by the
requested $p(C_3)=0.0,\ldots,1.0$ \cite{perera2020}.

\begin{table*}[t]
\caption{Benchmark design.  A \qefem{} trial contains the full replica
population.  G-set reports accumulated adaptive search and frozen evaluation
separately; its retained best is not interpreted as a fixed-trial success
probability.}
\label{tab:benchmark-design}
\centering
\small
\setlength{\tabcolsep}{5.5pt}
\begin{tabular}{lrrrrll}
\toprule
Family & instances & spins & structure & trials/instance & metric & comparator \\
\midrule
G-set G1--G54 & 54 & 800--3000 & sparse, varied & adaptive $+$ 2--5 eval. & cut and reference gap & -- \\
K2000 & 1 & 2000 & complete signed & 5 & cut and reference gap & \lqa{} \\
Wishart & 10 & 500 & dense planted & 5 & relative energy error, hits & \lqa{} \\
Chook & 11 & 1024 & periodic planted & 10 & relative energy error, hits & \lqa{} \\
\bottomrule
\end{tabular}
\end{table*}

For MaxCut, the residual is $C_{\rm ref}-C$, where $C_{\rm ref}$ is the stored
benchmark target.  These targets are quality references, not proofs of global
optimality.  For planted instances, where the ground energy $E_0$ is known,
we report $\epsilon=|E-E_0|/|E_0|$ and the number of trials that reach $E_0$.
Every final spin population is independently rescored against the original
instance.

\subsection{Frozen comparison settings}

The comparison uses the strongest frozen configurations available for each
method rather than forcing both methods into an artificial common optimizer.
This measures solution quality under the stated method-native budgets.  It is
not an equal-time or equal-operation comparison: \qefem{} evolves 8192 product
states per trial, whereas \lqa{} evolves one latent state.

\begin{table*}[t]
\caption{Frozen \qefem{} settings.  The RMSprop tuple is
$(\alpha_{\rm RMS},m,\mathrm{wd})$; $f_\Gamma$ is the fraction of the run at
which the transverse field reaches its terminal value.  Temperature and
transverse schedules are linear.}
\label{tab:qefem-settings}
\centering
\small
\setlength{\tabcolsep}{4.2pt}
\begin{tabular}{lrrrrclllc}
\toprule
Family & steps & replicas & trials & $\eta$ & RMSprop $(\alpha,m,\mathrm{wd})$ & $h$ & $T_{\max}\to T_{\min}$ & $\Gamma_{\max}\to\Gamma_{\min}$ & $f_\Gamma$ \\
\midrule
K2000  & 1000 & 8192 & 5  & .03 & $(.98,.91,.010)$ & .10 & $1\to10^{-4}$ & $1\to0$ & .50 \\
Wishart& 500  & 8192 & 5  & .03 & $(.56,.63,.013)$ & .10 & $1\to10^{-4}$ & $.3\to0$ & .50 \\
Chook  & 500  & 8192 & 10 & .03 & $(.90,.90,.001)$ & .05 & $2\to10^{-3}$ & $1\to0$ & 1.00 \\
\bottomrule
\end{tabular}
\end{table*}

The couplings are divided by their root-mean-square scale during optimization
and restored for scoring.  The corresponding scales are 44.71018 for K2000,
0.316598 for Wishart, and 3.162278 for Chook.  All three use dense float32
tensors.  The five K2000 and Wishart trials use one frozen seed batch; Chook
combines two batches of five trials.

\begin{table}[t]
\caption{Frozen \lqa{} settings.  Adam uses
$(\beta_1,\beta_2)=(.9,.999)$, zero weight decay, and one state.}
\label{tab:lqa-settings}
\centering
\small
\setlength{\tabcolsep}{4pt}
\begin{tabular}{lrrrr}
\toprule
Family & steps & $\eta$ & $g$ & $f$ \\
\midrule
K2000 & 5000 & 1 & .1 & .1 \\
Wishart & 500 & 2 & 1 & .1 \\
Chook & 500 & 2 & 1 & .1 \\
\bottomrule
\end{tabular}
\end{table}

\subsection{G-set search and confirmation}

The G-set evidence was accumulated in two stages.  The first stage evaluated
the smooth objective across G1--G54 and reached 23 references.  The second
stage revisited the 31 unmatched graphs with spectral conditioning and manual
exact or discrete-neighbour updates.  Candidate configurations were promoted
and then evaluated with new seed bases.  This distinction matters: the best
cut found during adaptive search is evidence of attainable solution quality,
whereas fresh-seed recovery measures whether a frozen configuration reproduces
that attainment.

\begin{table}[t]
\caption{G-set protocol ranges.  Exact graph-level settings and seeds are
retained in the reproducibility record.}
\label{tab:gset-protocol}
\centering
\small
\setlength{\tabcolsep}{4pt}
\begin{tabular}{lrrrr}
\toprule
Stage & graphs & steps & replicas & eval. trials \\
\midrule
Smooth & 54 & 1000--8000 & 4096 & 5 \\
Revisited & 31 & 4000--8000 & 2048--4096 & 2--3 \\
\bottomrule
\end{tabular}
\end{table}

The smooth stage uses RMS coupling normalization, $T_{\min}=10^{-4}$,
$\Gamma_{\max}=1$, and either linear or quadratic cooling.  The revisited
stage uses graph-specific spectral schedules, $\Gamma\to0$ over the full run,
SGD momentum .95, and weight decay .01.  The broad learning-rate and schedule
ranges reflect instance-specific tuning rather than a claim of a universal
G-set configuration.

\subsection{Reproducibility discipline}

Each endpoint is tied to an input hash, code revision, configuration, seed,
device, and numerical precision.  Final binary populations are rescored
independently, and a stored trajectory follows one coherent replica selected
by final readout.  The K2000, Wishart, and Chook tables below are frozen
evaluations rather than values selected after inspecting the test trials.  A
separate CPU replay of one K2000 seed reproduces its archived cut exactly,
although identical trajectories are not expected across software versions
and devices when the initialized fields differ at the bit level.

These rules separate adaptive attainment from frozen evaluation and solution
quality from timing.  The results below follow the same order as the protocol:
G-set first, then the dense K2000 instance, and finally the two planted
families.

\section{Results}
\label{sec:results}

\subsection{Overview}

Table~\ref{tab:results-overview} gives the aggregate comparison before the
instance-level results.  For the planted suites, every instance receives equal
weight in the reported mean; the following subsections then expose the
instance- and trial-level evidence behind those aggregates.

\begin{table}[h]
\caption{Aggregate frozen results.  ``Best'' is the largest cut or smallest
relative error; hits count trials exactly at the planted energy.}
\label{tab:results-overview}
\centering
\scriptsize
\setlength{\tabcolsep}{3pt}
\begin{tabular}{llrr}
\toprule
Family & method (trials) & mean & best; hits \\
\midrule
K2000 & \qefem{} (5) & cut 33281.2 & cut 33293; gap 44 \\
K2000 & \lqa{} (5) & cut 33249.0 & cut 33249; gap 88 \\
Wishart & \qefem{} (50) & $\epsilon=4.4783\%$ & $0$; $5/50$ \\
Wishart & \lqa{} (50) & $\epsilon=6.9767\%$ & $2.1347\%$; $0/50$ \\
Chook & \qefem{} (110) & $\epsilon=.21985\%$ & $0$; $60/110$ \\
Chook & \lqa{} (110) & $\epsilon=1.11142\%$ & $0$; $10/110$ \\
\bottomrule
\end{tabular}
\end{table}

\subsection{G-set target attainment}

The accumulated G1--G54 study reaches 33 stored reference cuts without local
postprocessing.  The smooth stage accounts for 23 matches.  Revisiting only
the 31 unmatched graphs adds ten: G15--G19, G21, G22, G24, G25, and G28.
Eight of these ten are recovered after freezing the configuration and changing
the seeds.  G17 and G28 remain search-stage attainments rather than fresh-seed
recoveries.  Across the 31 revisited graphs, 27 retained cuts improve over the
preceding best.

\begin{table}[t]
\caption{Accumulated G1--G54 evidence.  Fresh recovery is measured only after
configuration selection and therefore is not a success rate over the adaptive
search.}
\label{tab:gset-summary}
\centering
\small
\setlength{\tabcolsep}{5pt}
\begin{tabular}{lr}
\toprule
Evidence & matches \\
\midrule
Smooth study & $23/54$ \\
Additional revisited matches & $10$ \\
Fresh recovery of new matches & $8/10$ \\
Combined retained & $33/54$ \\
\bottomrule
\end{tabular}
\end{table}

The unresolved cases are not uniform failures.  G14, G29, G30, G31, and G51
are each one cut unit below their stored references; G23, G26, and G40 are two
units below.  The largest residuals in this group occur on G35--G38, where the
remaining gaps are 13--18.  Reporting the residual together with the attained
cut is essential: a one-unit miss on a large coupled graph demonstrates a very
different level of solution quality from a distant endpoint.  The complete
54-instance table is given in Appendix~\ref{app:gset-results}.

\subsection{Dense K2000 graph}

K2000 tests the method in a qualitatively different regime: every vertex is
connected to every other vertex.  The frozen \qefem{} trials produce cuts
between 33,274 and 33,293, with mean 33,281.2.  The best residual is 44, or
0.132\% of the stored target.  \lqa{} returns 33,249 in each of its five
stored trials, leaving residual 88.  Under these method-native settings, the
best \qefem{} endpoint therefore halves the unresolved target gap.

\begin{table}[t]
\caption{Every frozen K2000 trial.  Both methods use the same portable seed in
each row, although their internal state dimensions differ.}
\label{tab:k2000-trials}
\centering
\scriptsize
\setlength{\tabcolsep}{3.2pt}
\begin{tabular}{rrrrr}
\toprule
seed & \qefem{} cut & gap & \lqa{} cut & gap \\
\midrule
3371835386 & 33274 & 63 & 33249 & 88 \\
3700479728 & 33274 & 63 & 33249 & 88 \\
1934420243 & 33274 & 63 & 33249 & 88 \\
716463040  & 33293 & 44 & 33249 & 88 \\
424211264  & 33291 & 46 & 33249 & 88 \\
\midrule
mean & 33281.2 & 55.8 & 33249.0 & 88.0 \\
\bottomrule
\end{tabular}
\end{table}

The gain should not be misread as a timing comparison.  \qefem{} uses a large
batched population, and the present experiment was designed to test attainable
quality rather than matched computational work.  What the result establishes
is that the mixed-state relaxation remains effective even when the objective
contains nearly two million pairwise interactions.

\subsection{Wishart planted ensemble}

The planted benchmarks next test whether the same relaxation remains effective
when the optimum is known rather than represented by a stored reference.
Wishart varies the ruggedness of a dense planted problem through $\alpha$.
\qefem{} has lower mean relative error than \lqa{} at every tested value.  The
equal-instance mean falls from 6.9767\% to 4.4783\%, a 35.8\% reduction.  At
$\alpha=1$, all five \qefem{} trials reach the planted energy, whereas \lqa{}
retains mean absolute error 31.40.

\begin{table*}[t]
\caption{Wishart results, five trials per method and instance.  $\Delta E$ and
$\epsilon$ are trial means; hits count exact planted-energy endpoints.}
\label{tab:wishart-results}
\centering
\small
\setlength{\tabcolsep}{5pt}
\begin{tabular}{crrrrrr}
\toprule
$\alpha$ & \multicolumn{2}{c}{mean $\Delta E$} & \multicolumn{2}{c}{mean $\epsilon$ (\%)} & \multicolumn{2}{c}{hits/5} \\
\cmidrule(lr){2-3}\cmidrule(lr){4-5}\cmidrule(lr){6-7}
& \qefem{} & \lqa{} & \qefem{} & \lqa{} & \qefem{} & \lqa{} \\
\midrule
.1&.18&.64&.721&2.549&0&0\\
.2&.81&1.46&1.616&2.930&0&0\\
.3&1.97&2.96&2.617&3.936&0&0\\
.4&3.71&4.83&3.693&4.811&0&0\\
.5&6.02&7.71&4.808&6.155&0&0\\
.6&8.86&10.99&5.908&7.324&0&0\\
.7&12.57&14.78&7.192&8.456&0&0\\
.8&16.87&19.50&8.463&9.780&0&0\\
.9&21.94&25.32&9.766&11.272&0&0\\
1.0&0&31.40&0&12.551&5&0\\
\midrule
equal-instance mean&&&4.478&6.977&5/50&0/50\\
\bottomrule
\end{tabular}
\end{table*}

\subsection{Chook tile-planted ensemble}

The second planted family, Chook, shows a larger aggregate separation.
\qefem{} lowers the
equal-instance relative error from 1.11142\% to 0.21985\%, an 80.2\% reduction,
and improves the mean error at every requested $p(C_3)$.  For
$p(C_3)=0.5$ through 1.0, all 60 \qefem{} trials reach the planted energy.
The best \lqa{} behavior also occurs at the high-$p(C_3)$ end, but only ten of
110 trials reach the planted energy.

\begin{table*}[t]
\caption{Chook results, ten trials per method and instance.  The requested
$p(C_3)$ labels the generated instance; its realized tile fraction can differ
slightly.}
\label{tab:chook-results}
\centering
\small
\setlength{\tabcolsep}{5pt}
\begin{tabular}{crrrrrr}
\toprule
$p(C_3)$ & \multicolumn{2}{c}{mean $\Delta E$} & \multicolumn{2}{c}{mean $\epsilon$ (\%)} & \multicolumn{2}{c}{hits/10} \\
\cmidrule(lr){2-3}\cmidrule(lr){4-5}\cmidrule(lr){6-7}
& \qefem{} & \lqa{} & \qefem{} & \lqa{} & \qefem{} & \lqa{} \\
\midrule
.0&14.40&57.20&.703&2.793&0&0\\
.1&11.80&47.20&.594&2.374&0&0\\
.2&10.20&31.80&.524&1.635&0&0\\
.3&7.40&25.40&.389&1.337&0&0\\
.4&3.80&21.80&.208&1.192&0&0\\
.5&0&18.60&0&1.039&10&0\\
.6&0&15.80&0&.919&10&0\\
.7&0&7.20&0&.429&10&0\\
.8&0&5.20&0&.318&10&0\\
.9&0&2.20&0&.138&10&4\\
1.0&0&.80&0&.052&10&6\\
\midrule
equal-instance mean&&&.2199&1.1114&60/110&10/110\\
\bottomrule
\end{tabular}
\end{table*}

The Chook schedule retains the transverse field until the end of the run,
whereas K2000 and Wishart remove it halfway.  This family dependence is
important: the schedule is part of the practical solver and should be reported
explicitly rather than treated as incidental implementation detail.

\clearpage
\section{Discussion}
\label{sec:discussion}

The experiments support a simple conclusion: a factorized state need not be a
weak optimization model.  It cannot represent intersite correlations or
entanglement, yet its free-energy geometry can guide a large population toward
strong discrete states.  The transverse component is useful not because the
classical computer becomes quantum, but because it enlarges the continuous
path by which a longitudinal decision can form.  Entropy prevents premature
polarization early in the run; annealing then concentrates the longitudinal
components before sign readout produces a binary assignment.

Population search is equally central.  \qefem{} is not one deterministic
trajectory but thousands of independently initialized product states evolved
under the same objective.  The batch exposes parallel work and samples
multiple basins of the nonconvex landscape.  This improves attainable quality
at the cost of memory and total arithmetic.  The present comparisons therefore
support a quality claim, not a speedup claim.

The results also show why residuals must be interpreted rather than merely
plotted.  On K2000, gap 44 corresponds to 99.868\% of a strong stored target
on a complete graph.  On G-set, a residual of one means that the returned
binary state differs by only one unit of cut value from the reference.  These
near-target states demonstrate the capability of the method even when the
reference is not exactly reached.  Conversely, the larger residuals on
G35--G38 identify a genuine weakness and should remain visible.

Several limitations remain.  The G-set result combines adaptive search with
frozen confirmation and therefore does not estimate a single universal
success probability.  The \qefem{}--\lqa{} comparison uses method-native
resources rather than matched compute.  Hyperparameters are family dependent,
and sign readout can magnify small continuous differences near zero.  A larger
comparison should include controlled population scaling, time-to-target
curves, held-out instance families, and additional classical baselines.

\section{Conclusion}

\qefem{} turns a binary Ising objective into a product-state free energy,
explores that objective with many replicas, and returns an ordinary feasible
spin assignment.  The derivation is compact but complete: the Gibbs target
defines the variational principle, the product ansatz makes it tractable, and
the nonlinear Bloch fields provide a stable unconstrained parameterization.
The current evidence reaches 33 of 54 G-set references, comes within 44 of the
dense K2000 target, and lowers Wishart and Chook error relative to \lqa{} under
the stated settings.  The result is a classical solver with a simple state
model and a surprisingly strong reach across distinct large-scale binary
landscapes.

\section*{Author contributions}

R.D. conceived and developed the core \qefem{} method, implemented the
original solver, designed the analysis workflow, and drafted the manuscript.
F.P. contributed the nonlinear reparameterization, optimizer and schedule
refinement, and benchmark testing.  Both authors interpreted the results and
take responsibility for the final manuscript.

\begin{thebibliography}{99}

\bibitem{karp1972}
R. M. Karp, Reducibility among combinatorial problems, in
\emph{Complexity of Computer Computations}, edited by R. E. Miller and
J. W. Thatcher (Plenum, New York, 1972), pp. 85--103.

\bibitem{barahona1982}
F. Barahona, On the computational complexity of Ising spin glass models,
J. Phys. A: Math. Gen. \textbf{15}, 3241 (1982).

\bibitem{lucas2014}
A. Lucas, Ising formulations of many NP problems, Front. Phys. \textbf{2}, 5
(2014).

\bibitem{kirkpatrick1983}
S. Kirkpatrick, C. D. Gelatt, Jr., and M. P. Vecchi, Optimization by simulated
annealing, Science \textbf{220}, 671 (1983).

\bibitem{inagaki2016science}
T. Inagaki \emph{et al.}, A coherent Ising machine for 2000-node optimization
problems, Science \textbf{354}, 603 (2016).

\bibitem{goto2019sbm}
H. Goto, K. Tatsumura, and A. R. Dixon, Combinatorial optimization by
simulating adiabatic bifurcations in nonlinear Hamiltonian systems, Sci. Adv.
\textbf{5}, eaav2372 (2019).

\bibitem{aramon2019}
M. Aramon, G. Rosenberg, E. Valiante, T. Miyazawa, H. Tamura, and
H. G. Katzgraber, Physics-inspired optimization for quadratic unconstrained
problems using a digital annealer, Front. Phys. \textbf{7}, 48 (2019).

\bibitem{shen2025}
Z.-S. Shen \emph{et al.}, Free-energy machine for combinatorial optimization,
Nat. Comput. Sci. \textbf{5}, 322 (2025).

\bibitem{bowles2022}
J. Bowles, A. Dauphin, P. Huembeli, J. Martinez, and A. Ac\'in, Quadratic
unconstrained binary optimisation via quantum-inspired annealing, Phys. Rev.
Applied \textbf{18}, 034016 (2022).

\bibitem{wainwright2008}
M. J. Wainwright and M. I. Jordan, Graphical models, exponential families,
and variational inference, Found. Trends Mach. Learn. \textbf{1}, 1 (2008).

\bibitem{blei2017}
D. M. Blei, A. Kucukelbir, and J. D. McAuliffe, Variational inference: A
review for statisticians, J. Am. Stat. Assoc. \textbf{112}, 859 (2017).

\bibitem{nielsen2010}
M. A. Nielsen and I. L. Chuang, \emph{Quantum Computation and Quantum
Information}, 10th anniversary ed. (Cambridge University Press, Cambridge,
2010).

\bibitem{yeGset}
Y. Ye, G-set test problems, Stanford University,
\url{https://web.stanford.edu/~yyye/yyye/Gset/} (accessed September 2026).

\bibitem{hamze2020}
F. Hamze, J. Raymond, C. A. Pattison, K. Biswas, and H. G. Katzgraber,
Wishart planted ensemble: A tunably rugged pairwise Ising model with a
first-order phase transition, Phys. Rev. E \textbf{101}, 052102 (2020).

\bibitem{perera2020}
D. Perera, I. Akpabio, F. Hamze, S. Mandra, N. Rose, M. Aramon, and
H. G. Katzgraber, Chook---A comprehensive suite for generating binary
optimization problems with planted solutions, arXiv:2005.14344 (2020).

\bibitem{paszke2019}
A. Paszke \emph{et al.}, PyTorch: An imperative style, high-performance deep
learning library, in \emph{Advances in Neural Information Processing Systems
32} (2019), pp. 8024--8035.

\bibitem{tieleman2012}
T. Tieleman and G. Hinton, Lecture 6.5---RMSProp: Divide the gradient by a
running average of its recent magnitude, COURSERA: Neural Networks for Machine
Learning (2012).

\end{thebibliography}

\appendix

\newpage
\section{Complete G1--G54 endpoint table}
\label{app:gset-results}

\noindent
The table below lists every retained G-set endpoint.  Status O denotes an
endpoint attained in the original smooth study; F, attained again in frozen
fresh-seed evaluation; S, attained during adaptive search but not fresh
evaluation; and U, unresolved.

\begin{center}
\scriptsize
\setlength{\tabcolsep}{4.2pt}
\resizebox{\columnwidth}{!}{%
\begin{tabular}{lrrrc@{\hspace{10pt}}lrrrc}
\toprule
Graph & target & best & gap & status & Graph & target & best & gap & status \\
\midrule
G1&11624&11624&0&O&G28&3298&3298&0&S\\
G2&11620&11620&0&O&G29&3405&3404&1&U\\
G3&11622&11622&0&O&G30&3413&3412&1&U\\
G4&11646&11646&0&O&G31&3310&3309&1&U\\
G5&11631&11631&0&O&G32&1410&1406&4&U\\
G6&2178&2178&0&O&G33&1382&1376&6&U\\
G7&2006&2006&0&O&G34&1384&1378&6&U\\
G8&2005&2005&0&O&G35&7687&7674&13&U\\
G9&2054&2054&0&O&G36&7680&7664&16&U\\
G10&2000&2000&0&O&G37&7691&7673&18&U\\
G11&564&564&0&O&G38&7688&7674&14&U\\
G12&556&556&0&O&G39&2408&2404&4&U\\
G13&582&582&0&O&G40&2400&2398&2&U\\
G14&3064&3063&1&U&G41&2405&2402&3&U\\
G15&3050&3050&0&F&G42&2481&2473&8&U\\
G16&3052&3052&0&F&G43&6660&6660&0&O\\
G17&3047&3047&0&S&G44&6650&6650&0&O\\
G18&992&992&0&F&G45&6654&6654&0&O\\
G19&906&906&0&F&G46&6649&6649&0&O\\
G20&941&941&0&O&G47&6657&6657&0&O\\
G21&931&931&0&F&G48&6000&6000&0&O\\
G22&13359&13359&0&F&G49&6000&6000&0&O\\
G23&13344&13342&2&U&G50&5880&5880&0&O\\
G24&13337&13337&0&F&G51&3848&3847&1&U\\
G25&13340&13340&0&F&G52&3851&3848&3&U\\
G26&13328&13326&2&U&G53&3850&3846&4&U\\
G27&3341&3341&0&O&G54&3852&3848&4&U\\
\bottomrule
\end{tabular}}
\end{center}

\newpage
\section{QUBO conversion and quadratization}
\label{app:qubo}

For a symmetric matrix $Q$ stored in full, substitute
$\bm x=(\bm 1-\bm s)/2$ into $E_Q=\bm x^{\mathsf T}Q\bm x+c$:
\begin{align}
 E_Q(\bm s)
 &=\frac14(\bm 1-\bm s)^{\mathsf T}Q(\bm 1-\bm s)+c\nonumber\\
 &=c+\frac14\bm 1^{\mathsf T}Q\bm 1
   -\frac12(Q\bm 1)^{\mathsf T}\bm s
   +\frac14\bm s^{\mathsf T}Q\bm s.
 \label{eq:qubo-expansion}
\end{align}
Using $s_i^2=1$ and counting each unordered pair once gives
\begin{align}
 c_0&=c+\frac14\bm 1^{\mathsf T}Q\bm 1+\frac14\Tr Q,\nonumber\\
 h_i&=-\frac12(Q\bm 1)_i,
 &K_{ij}&=\frac12Q_{ij}.
 \label{eq:qubo-map}
\end{align}
Other storage conventions for off-diagonal QUBO terms produce equivalent
coefficients.  The implementation fixes one convention and verifies the
energy of converted assignments.

Higher-order binary products can be reduced with auxiliary variables.  To
replace $z=xy$, one convenient penalty is
\begin{equation}
 P\,[xy-2xz-2yz+3z], \qquad P>0,
 \label{eq:quadratization}
\end{equation}
which is zero exactly when $z=xy$ and positive on the other binary
assignments.  Choosing $P$ larger than the possible objective gain from
violating the constraint preserves the optimum.  Repeated substitutions
quadratize a higher-order polynomial at the cost of additional variables and
a wider coefficient range.

\section{Parameterization and analytic derivatives}
\label{app:param}

Write $\bm a_i=(u_i,v_i)$, $r_i=\|\bm a_i\|$, and
$g(r)=\tanh(r)/r$.  Then $\bm m_i=g(r_i)\bm a_i$.  The continuous extension
$g(0)=1$ removes the apparent singularity.  Its Jacobian is
\begin{equation}
 \frac{\partial\bm m}{\partial\bm a}
 =g(r)I+\frac{g'(r)}{r}\bm a\bm a^{\mathsf T},
 \qquad
 g'(r)=\frac{r\,\mathrm{sech}^2r-\tanh r}{r^2}.
 \label{eq:jacobian}
\end{equation}

The entropy in Eq.~\eqref{eq:local-entropy} may be evaluated stably in field
coordinates as
\begin{equation}
 S_i=\log(2\cosh r_i)-r_i\tanh r_i.
 \label{eq:entropy-field}
\end{equation}
For the general Ising objective, derivatives in physical coordinates are
\begin{align}
 \frac{\partial\mathcal F}{\partial m_{z,i}}
 &=h_i+\sum_{j\ne i}K_{ij}m_{z,j}
   +T\,\operatorname{arctanh}(R_i)\frac{m_{z,i}}{R_i},\\
 \frac{\partial\mathcal F}{\partial m_{x,i}}
 &=-\Gamma
   +T\,\operatorname{arctanh}(R_i)\frac{m_{x,i}}{R_i}.
 \label{eq:physical-gradient}
\end{align}
Multiplication by Eq.~\eqref{eq:jacobian} gives derivatives with respect to
$(u_i,v_i)$.  The limits at $R_i=0$ are evaluated analytically.  These
expressions agree with automatic differentiation of
Eq.~\eqref{eq:qefem-objective}.

\section{Schedules and discrete-neighbour continuation}
\label{app:schedules}

For normalized progress $\tau=t/(L-1)$, the schedule family is
\begin{equation}
 y_t=y_{\min}+(y_{\max}-y_{\min})
 \left[1-\min\!\left(\frac{\tau}{f},1\right)\right]^p.
 \label{eq:schedule}
\end{equation}
Here $y$ denotes $T$ or $\Gamma$, $f$ is the fraction at which the terminal
value is reached, and $p=1$ gives a linear schedule.  The G-set smooth stage
also uses $p=2$ cooling.  Setting $f_{\Gamma}=1/2$ removes the driver halfway
through the run while leaving a classical optimization tail.

The exact longitudinal interaction is $(K\bm m_z)_i$.  The G-set continuation
replaces the neighbour vector by
\begin{equation}
 \widetilde{\bm m}_z^{(t)}
 =(1-\lambda_t)\bm m_z^{(t)}
 +\lambda_t\sign(\bm m_z^{(t)}).
 \label{eq:discrete-neighbour}
\end{equation}
Setting $\lambda_t=1$ replaces the neighbour magnetization by its sign
throughout the active region; alternatively, $\lambda_t$ can increase
gradually.  Entropy, transverse, and field-map derivatives remain
unchanged.  Because this continuation changes only the search direction, the
final binary state is still audited against the original graph.

\end{document}
